\section{Intuition}

\emph{Note: This originally was part of the notes for an open class in non-duality, given in the summer of 2005.}

In \textbf{The Act of Creation}, Arthur Koestler writes:

\begin{quotex}
Not only scientists, painters, and musicians find it often difficult to convert their ideas into verbal currency, but writers too. Even H. G. Wells lamented: ``The forceps of our minds are clumsy things and crush the truth a little in the course of taking hold of it.'' The novelist suffers from the poverty of his vocabulary when he tries to describe what his characters feel (as distinct from what they think or do). He can write streams of what goes on in the cranial cavity, but if it is a pain in the abdominal cavity, all he can say is `it hurts' -- or use some equally insipid synonym. Suffering is `dumb'; the glandular and visceral processes which colour emotions do not lend themselves to verbal articulation.
\end{quotex}


In our verbal, hyper-rationalistic age, we privilege `thinking' and `doing', and consider `feeling' to be merely subjective and therefore somehow less real. The opposite is actually the case. Our `feelings' give us direct access to the inner workings of the cosmos, and our `thinkings' and `doings' are superimpositions.

Of course, to the writer Wells' chagrin, our feelings can only be known directly, in consciousness --- this direct knowing is precisely what is being meant by the term ``intuition''. Basically, then, to know an emotion such as anger, or joy, or love, we must be angry, joyous, or loving; no verbal description or brain scan can replace that direct acquaintance. Furthermore, when we learn to experience the world as arising continuously in consciousness, we will then know the world directly, by intuition. And when we then come to know the Self that is the unchanging presence in every act of consciousness, we will be realized. It is that simple.


\flrightit{Posted on 2010-10-24 by Cologero}

